
% BEGIN bib
 
\par\medskip\noindent\textbf{References}\par
  
\par\noindent  [1]  Majid Alizadeh, Mohammad Ardeshir, Philippe Balbiani, and Mojtaba Mojtahedi.  Unification types in Euclidean modal logics.  Log. J. IGPL, 31(3):422–440, 2023.  doi:10.1093/JIGPAL/JZAB036.  
 
\par\noindent  [2]  Franz Baader.  The theory of idempotent semigroups is of unification type zero.  J. Autom. Reason., 2(3):283–286, 1986.  doi:10.1007/BF02328451.  
 
\par\noindent  [3]  Franz Baader.  Characterization of unification type zero.  In Nachum Dershowitz, editor, Rewriting Techniques and Applications, 3rd International Conference, RTA-89, Proceedings, volume 355 of Lecture Notes in Computer Science, pages 2–14. Springer, 1989.  doi:10.1007/3-540-51081-8\_96.  
 
\par\noindent  [4]  Franz Baader.  Unification, weak unification, upper bound, lower bound, and generalization problems.  In Ronald V. Book, editor, Rewriting Techniques and Applications, 4th International Conference, RTA-91, Proceedings, volume 488 of Lecture Notes in Computer Science, pages 86–97. Springer, 1991.  doi:10.1007/3-540-53904-2\_88.  
 
\par\noindent  [5]  Franz Baader.  Unification in commutative theories, Hilbert’s basis theorem, and Gröbner bases.  J. ACM, 40(3):477–503, 1993.  doi:10.1145/174130.174133.  
 
\par\noindent  [6]  Franz Baader, Sebastian Brandt, and Carsten Lutz.  Pushing the $\mathcal{EL}$ envelope.  In Leslie Pack Kaelbling and Alessandro Saffiotti, editors, IJCAI-05, Proceedings of the Nineteenth International Joint Conference on Artificial Intelligence, pages 364–369. Professional Book Center, 2005.  URL: http://ijcai.org/Proceedings/05/Papers/0372.pdf.  
 
\par\noindent  [7]  Franz Baader and Wolfram Büttner.  Unification in commutative idempotent monoids.  Theor. Comput. Sci., 56:345–353, 1988.  doi:10.1016/0304-3975(88)90140-5.  
 
\par\noindent  [8]  Franz Baader and Silvio Ghilardi.  Unification in modal and description logics.  Log. J. IGPL, 19(6):705–730, 2011.  doi:10.1093/JIGPAL/JZQ008.  
 
\par\noindent  [9]  Franz Baader and Pierre Ludmann.  The exact unification type of commutative theories.  In Santiago Escobar and Mateu Villaret, editors, Proceedings of the 29th International Workshop on Unification (UNIF’15), pages 19–23, 2015.  URL: https://lat.inf.tu-dresden.de/research/papers/2015/BaLu-UNIF2015.pdf.  
 
\par\noindent  [10]  Franz Baader and Pierre Ludmann.  The unification type of ACUI w.r.t. the unrestricted instantiation preorder is not finitary.  In Silvio Ghilardi and Manfred Schmidt-Schauß, editors, Proceedings of the 30th International Workshop on Unification, UNIF 2016, pages 31–36, 2016.  URL: http://users.mat.unimi.it/users/ghilardi/UNIF2016/UNIF16-abstracts.pdf\#page=31.  
 
\par\noindent  [11]  Franz Baader, Carsten Lutz, and Boontawee Suntisrivaraporn.  CEL - A polynomial-time reasoner for life science ontologies.  In Ulrich Furbach and Natarajan Shankar, editors, Automated Reasoning, Third International Joint Conference, IJCAR 2006, Proceedings, volume 4130 of Lecture Notes in Computer Science, pages 287–291. Springer, 2006.  doi:10.1007/11814771\_25.  
 
\par\noindent  [12]  Franz Baader and Barbara Morawska.  Unification in the description logic $\mathcal{EL}$.  In Ralf Treinen, editor, Rewriting Techniques and Applications, 20th International Conference, RTA 2009, Proceedings, volume 5595 of Lecture Notes in Computer Science, pages 350–364. Springer, 2009.  doi:10.1007/978-3-642-02348-4\_25.  
 
\par\noindent  [13]  Franz Baader and Barbara Morawska.  Unification in the description logic $\mathcal{EL}$.  Log. Methods Comput. Sci., 6(3), 2010.  URL: http://arxiv.org/abs/1006.2289.  
 
\par\noindent  [14]  Franz Baader and Paliath Narendran.  Unification of concept terms in description logics.  J. Symb. Comput., 31(3):277–305, 2001.  doi:10.1006/JSCO.2000.0426.  
 
\par\noindent  [15]  Franz Baader, Thanh Binh Nguyen, Stefan Borgwardt, and Barbara Morawska.  Deciding unifiability and computing local unifiers in the description logic $\mathcal{EL}$ without top constructor.  Notre Dame J. Formal Log., 57(4):443–476, 2016.  doi:10.1215/00294527-3555507.  
 
\par\noindent  [16]  Franz Baader and Tobias Nipkow.  Term Rewriting and All That.  Cambridge University Press, 1998.  doi:10.1017/CBO9781139172752.  
 
\par\noindent  [17]  Franz Baader and Werner Nutt.  Combination problems for commutative/monoidal theories or how algebra can help in equational unification.  Appl. Algebra Eng. Commun. Comput., 7(4):309–337, 1996.  doi:10.1007/BF01195536.  
 
\par\noindent  [18]  Franz Baader and Jörg H. Siekmann.  Unification theory.  In Dov M. Gabbay, Christopher J. Hogger, J. A. Robinson, and Jörg H. Siekmann, editors, Handbook of Logic in Artificial Intelligence and Logic Programming, Volume 2, Deduction Methodologies, pages 41–126. Oxford University Press, 1994.  doi:10.1093/oso/9780198537465.003.0002.  
 
\par\noindent  [19]  Franz Baader and Wayne Snyder.  Unification theory.  In John Alan Robinson and Andrei Voronkov, editors, Handbook of Automated Reasoning, volume 1, pages 445–532. Elsevier and MIT Press, 2001.  doi:10.1016/B978-044450813-3/50010-2.  
 
\par\noindent  [20]  Leo Bachmair and Nachum Dershowitz.  Completion for rewriting modulo a congruence.  Theor. Comput. Sci., 67(2\&3):173–201, 1989.  doi:10.1016/0304-3975(89)90003-0.  
 
\par\noindent  [21]  Philippe Balbiani and Çiğdem Gencer.  About the unification types of modal logics.  In Nick Bezhanishvili, Rosalie Iemhoff, and Fan Yang, editors, Dick de Jongh on Intuitionistic and Provability Logics, pages 175–202. Springer International Publishing, 2024.  doi:10.1007/978-3-031-47921-2\_7.  
 
\par\noindent  [22]  Philippe Balbiani, Çigdem Gencer, Maryam Rostamigiv, and Tinko Tinchev.  Remarks about the unification types of some locally tabular normal modal logics.  Log. J. IGPL, 31(1):115–139, 2023.  doi:10.1093/JIGPAL/JZAB033.  
 
\par\noindent  [23]  Leonardo Manuel Cabrer and George Metcalfe.  From admissibility to a new hierarchy of unification types.  In Temur Kutsia and Christophe Ringeissen, editors, Proceedings of the 28th International Workshop on Unification, UNIF 2014, pages 41–46, 2014.  URL: https://www3.risc.jku.at/publications/download/risc\_5001/proceedings-UNIF2014.pdf\#page=46.  
 
\par\noindent  [24]  Leonardo Manuel Cabrer and George Metcalfe.  Exact unification and admissibility.  Log. Methods Comput. Sci., 11(3), 2015.  doi:10.2168/LMCS-11(3:23)2015.  
 
\par\noindent  [25]  Jacques Corbin and Michel Bidoit.  A rehabilitation of Robinson’s unification algorithm.  In R. E. A. Mason, editor, Proceedings of the 9th World Computer Congress, IFIP’83, pages 909–914, Paris, France, 1983. Elsevier.  
 
\par\noindent  [26]  Wojciech Dzik, Slawomir Kost, and Piotr Wojtylak.  Finitary unification in locally tabular modal logics characterized.  Ann. Pure Appl. Log., 173(4):103072, 2022.  doi:10.1016/J.APAL.2021.103072.  
 
\par\noindent  [27]  Wojciech Dzik, Slawomir Kost, and Piotr Wojtylak.  Unification types and union splittings in intermediate logics.  Ann. Pure Appl. Log., 176(1):103508, 2025.  doi:10.1016/J.APAL.2024.103508.  
 
\par\noindent  [28]  Elmar Eder.  Properties of substitutions and unifications.  J. Symb. Comput., 1(1):31–46, 1985.  doi:10.1016/S0747-7171(85)80027-4.  
 
\par\noindent  [29]  François Fages and Gérard P. Huet.  Complete sets of unifiers and matchers in equational theories.  In Giorgio Ausiello and Marco Protasi, editors, CAAP’83, Trees in Algebra and Programming, 8th Colloquium, Proceedings, volume 159 of Lecture Notes in Computer Science, pages 205–220. Springer, 1983.  doi:10.1007/3-540-12727-5\_12.  
 
\par\noindent  [30]  François Fages and Gérard P. Huet.  Complete sets of unifiers and matchers in equational theories.  Theor. Comput. Sci., 43:189–200, 1986.  doi:10.1016/0304-3975(86)90175-1.  
 
\par\noindent  [31]  Silvio Ghilardi.  Unification in intuitionistic logic.  J. Symb. Log., 64(2):859–880, 1999.  doi:10.2307/2586506.  
 
\par\noindent  [32]  Miki Hermann, Laurent Juban, and Phokion G. Kolaitis.  On the complexity of counting the Hilbert basis of a linear diophnatine system.  In Harald Ganzinger, David A. McAllester, and Andrei Voronkov, editors, Logic Programming and Automated Reasoning, 6th International Conference, LPAR’99, Proceedings, volume 1705 of Lecture Notes in Computer Science, pages 13–32. Springer, 1999.  doi:10.1007/3-540-48242-3\_2.  
 
\par\noindent  [33]  Michael Hoche, Jörg H. Siekmann, and Peter Szabó.  String unification is essentially infinitary.  FLAP, 3(5):755–788, 2016.  URL: http://www.collegepublications.co.uk/downloads/ifcolog00009.pdf.  
 
\par\noindent  [34]  Rosalie Iemhoff.  A syntactic approach to unification in transitive reflexive modal logics.  Notre Dame J. Formal Log., 57(2):233–247, 2016.  doi:10.1215/00294527-345997.  
 
\par\noindent  [35]  Emil Jerábek.  Blending margins: the modal logic K has nullary unification type.  J. Log. Comput., 25(5):1231–1240, 2015.  doi:10.1093/LOGCOM/EXT055.  
 
\par\noindent  [36]  Jean-Pierre Jouannaud and Claude Kirchner.  Solving equations in abstract algebras: A rule-based survey of unification.  In J.-L. Lassez and G. Plotkin, editors, Computational Logic: Essays in Honor of A. Robinson. MIT Press, Cambridge, MA, 1991.  
 
\par\noindent  [37]  Jean-Pierre Jouannaud and Helene Kirchner.  Completion of a set of rules modulo a set of equations.  SIAM J. Computing, 15(4):1155–1194, 1986.  doi:10.1137/0215084.  
 
\par\noindent  [38]  Donald E. Knuth and Peter B. Bendix.  Simple word problems in universal algebras.  In J. Leech, editor, Computational Problems in Abstract Algebra. Pergamon Press, Oxford, 1970.  doi:10.1016/B978-0-08-012975-4.50028-X.  
 
\par\noindent  [39]  Ralf Küsters.  Non-Standard Inferences in Description Logics, volume 2100 of Lecture Notes in Computer Science.  Springer, 2001.  doi:10.1007/3-540-44613-3.  
 
\par\noindent  [40]  Dallas Lankford.  Non-negative integer basis algorithms for linear equations with integer coefficients.  J. Autom. Reason., 5(1):25–35, 1989.  doi:10.1007/BF00245019.  
 
\par\noindent  [41]  Jean-Louis Lassez, Michael J. Maher, and Kim Marriott.  Unification revisited.  In Jack Minker, editor, Foundations of Deductive Databases and Logic Programming, pages 587–625. Morgan Kaufmann, 1988.  doi:10.1016/B978-0-934613-40-8.50019-1.  
 
\par\noindent  [42]  Alberto Martelli and Ugo Montanari.  An efficient unification algorithm.  ACM Trans. Program. Lang. Syst., 4(2):258–282, 1982.  doi:10.1145/357162.357169.  
 
\par\noindent  [43]  Mike Paterson and Mark N. Wegman.  Linear unification.  J. Comput. Syst. Sci., 16(2):158–167, 1978.  doi:10.1016/0022-0000(78)90043-0.  
 
\par\noindent  [44]  Gerald E. Peterson and Mark E. Stickel.  Complete sets of reductions for some equational theories.  J. of the ACM, 28(2):233–264, 1981.  doi:10.1145/322248.322251.  
 
\par\noindent  [45]  Gordon Plotkin.  Building in equational theories.  Machine Intelligence, 7:73–90, 1972.  
 
\par\noindent  [46]  Peter Raulefs, Jörg H. Siekmann, Peter Szabó, and E. Unvericht.  A short survey on the state of the art in matching and unification problems.  SIGSAM Bull., 13(2):14–20, 1979.  doi:10.1145/1089208.1089210.  
 
\par\noindent  [47]  John Alan Robinson.  A machine oriented logic based on the resolution principle.  J. of the ACM, 12(1):23–41, 1965.  doi:10.1145/321250.321253.  
 
\par\noindent  [48]  Manfred Schmidt-Schauß.  Unification under associativity and idempotence is of type nullary.  J. Autom. Reason., 2(3):277–281, 1986.  doi:10.1007/BF02328450.  
 
\par\noindent  [49]  Manfred Schmidt-Schauß.  A decision algorithm for distributive unification.  Theor. Comput. Sci., 208(1-2):111–148, 1998.  doi:10.1016/S0304-3975(98)00081-4.  
 
\par\noindent  [50]  Jörg H. Siekmann.  Matching under commutativity.  In Edward W. Ng, editor, Symbolic and Algebraic Computation, EUROSAM ’79, An International Symposium on Symbolic and Algebraic Computation, Proceedings, volume 72 of Lecture Notes in Computer Science, pages 531–545. Springer, 1979.  doi:10.1007/3-540-09519-5\_101.  
 
\par\noindent  [51]  Jörg H. Siekmann.  Universal unification.  In Robert E. Shostak, editor, 7th International Conference on Automated Deduction, CADE 7, Proceedings, volume 170 of Lecture Notes in Computer Science, pages 1–42. Springer, 1984.  doi:10.1007/978-0-387-34768-4\_1.  
 
\par\noindent  [52]  Jörg H. Siekmann.  Unification theory.  In Benedict du Boulay, David C. Hogg, and Luc Steels, editors, Advances in Artificial Intelligence II, Seventh European Conference on Artificial Intelligence, ECAI 1986, Proceedings, pages 365–400. North-Holland, 1986.  
 
\par\noindent  [53]  Jörg H. Siekmann.  Unification theory.  J. Symb. Comput., 7(3/4):207–274, 1989.  doi:10.1016/S0747-7171(89)80012-4.  
 
\par\noindent  [54]  Viorica Sofronie-Stokkermans.  Locality and applications to subsumption testing in $\mathcal{EL}$ and some of its extensions.  Sci. Ann. Comput. Sci., 23(2):251–284, 2013.  doi:10.7561/SACS.2013.2.251.  
 
\par\noindent  [55]  Viorica Sofronie-Stokkermans.  Representation theorems and locality for subsumption testing and interpolation in the description logics $\mathcal{EL}$, $\mathcal{EL}^{+}$ and their extensions with $n$-ary roles and numerical domains.  Fundam. Informaticae, 156(3-4):361–411, 2017.  doi:10.3233/FI-2017-1612.  
 
\par\noindent  [56]  Mark E. Stickel.  A unification algorithm for associative-commutative functions.  J. ACM, 28(3):423–434, 1981.  doi:10.1145/322261.322262.  
 
\par\noindent  [57]  Peter Szabó and Jörg H. Siekmann.  $E$-unification based on generalized embedding.  Math. Struct. Comput. Sci., 31(8):898–917, 2021.  doi:10.1017/S0960129522000019.  
 
\par\noindent  [58]  Katherine A. Yelick.  Combining unification algorithms for confined regular equational theories.  In Jean-Pierre Jouannaud, editor, Rewriting Techniques and Applications, First International Conference, RTA-85, Proceedings, volume 202 of Lecture Notes in Computer Science, pages 365–380. Springer, 1985.  doi:10.1007/3-540-15976-2\_18.  
  
% END bib
